\documentclass[12pt]{article} % 12pt article document

% ---- encoding and language ----
\usepackage[utf8]{inputenc} % UTF-8 input encoding
\usepackage[T2A]{fontenc}   % Cyrillic font encoding
\usepackage[english,russian]{babel} % languages: English + Russian

% ---- layout and appearance ----
\usepackage[a4paper, left=20mm, right=15mm, top=16mm, bottom=18mm]{geometry} % page margins
\usepackage{paracol}       % parallel columns
\usepackage{graphicx}      % images
\usepackage{xcolor}        % colors
\usepackage{setspace}      % line spacing
\usepackage{float}         % control float placement
\usepackage{microtype}     % better typographic spacing
\usepackage{array}         % better tables
\usepackage{ragged2e}      % ragged text alignment

\usepackage{tempora}       % Times-like font

\pagestyle{empty}          % no page numbers
\setlength{\columnsep}{14pt} % spacing between columns

\newcommand{\magentaTitle}[1]{\noindent\textcolor[rgb]{0.99,0.05,0.75}{\textbf{#1}}} % colored section title

\begin{document}

\columnratio{0.60,0.40} % ratio between the two columns

\begin{paracol}{2} % begin 2-column environment

% LEFT COLUMN
\begin{leftcolumn}

\begin{figure}[H] % image with fixed placement
  \centering
  \includegraphics[width=\linewidth]{image.png} % image fills column width
  {\small Рис. 7.\par} % small caption
\end{figure}

\vspace{6pt}

\magentaTitle{Шаг 2 -- изучаем удобные функции.} % section title in magenta

Если \textit{f} и \textit{g} -- удобные функции, а \textit{c} и \textit{d} -- действительные числа, то функция \textit{cf + dg}, которая по определению принимает на фигуре \textit{F} значение $cf(F)+dg(F)$, очевидно, является удобной.

\begin{spacing}{1.35} % increased line spacing
Итак, с удобными функциями можно обращаться как с векторами: их можно умножать на действительные числа и складывать. Положим, например, \textit{P = N\textsubscript{2} -- 2N + E} (т. е. $P(F)=N(2F)-2N(F)+E(F)$). Так как функции \textit{N\textsubscript{2}}, \textit{N} и \textit{E} -- удобные, то и функция \textit{P} -- удобная.
\end{spacing}

Сопоставим теперь каждой удобной функции $f\in Y$ вектор в трёхмерном пространстве, а именно,

% FIXED NON-OVERLAPPING FORMULA 
\begin{table}[H]
\centering
\small
\begin{tabular}{@{}c@{}}
$(f(X_1);\; f(X_2);\; f(X_3))$ % column vector notation
\end{tabular}
\end{table}


где $X_1, X_2, X_3$ -- три целочисленные фигуры, изображённые на рисунке~7. Этот вектор будем обозначать через $\vec f$. В первых строках таблицы 1 записано, какие векторы сопоставляются уже известным нам функциям.

\vspace{6pt}
\noindent\textbf{Т а б л и ц а 1} % table title

\small
\setlength{\tabcolsep}{4pt}    % cell padding
\renewcommand{\arraystretch}{1.02} % row height
\begin{center}
\begin{tabular}{|c|c|c|c|c|} % bordered table
  \hline
  & $f$ & $f(X_1)$ & $f(X_2)$ & $f(X_3)$ \\ \hline
  1 & $S$ & 0 & 0 & $1/2$ \\ \hline
  2 & $N$ & 1 & 2 & 3 \\ \hline
  3 & $E$ & 1 & 1 & 1 \\ \hline
  4 & $N_2$ & 1 & 3 & 6 \\ \hline
  5 & $P$ & 0 & 0 & 1 \\ \hline
  6 & ? & 0 & 2 & 3 \\ \hline
  7 & $e_1=2E+2S-N$ & 1 & 0 & 0 \\ \hline
  8 & $e_2=N-E-4S$ & 0 & 1 & 0 \\ \hline
  9 & $e_3=2S$ & 0 & 0 & 1 \\ \hline
\end{tabular}
\end{center}
\normalsize

\end{leftcolumn}

% RIGHT COLUMN 
\switchcolumn % move to second column

 % left alignment without justification

А всякий ли вектор соответствует какой-нибудь удобной функции (например, вектор $(0;2;3)$ из шестой строки)? Положительный ответ дают строки 7--9 таблицы 1. Из них видно, что удобным функциям $e_1,e_2,e_3$ соответствуют базисные векторы в трёхмерном пространстве. Поэтому для любых чисел $c_1,c_2,c_3$ можно указать удобную функцию $c_1e_1+c_2e_2+c_3e_3$.

\bigskip

\magentaTitle{Шаг 3 -- критерий совпадения удобных функций.} % new subsection

До сих пор мы нигде не пользовались тем, что имеем дело с удобными функциями. Любой функции на $\Phi$ можно сопоставить вектор трёхмерного пространства. Например, функции $R$ из упражнения 2 соответствует вектор из шестой строки таблицы 1. Особое «удобство» удобных функций заключается в том, что между ними и векторами трёхмерного пространства имеется \textit{взаимно однозначное соответствие}. Иными словами, если $f$ и $g$ -- удобные функции и $\vec f=\vec g$, то $f=g$.

Обозначим разность $f-g$ через $h$. Тогда утверждение о взаимной однозначности принимает такой вид.

\medskip

\noindent\textbf{Основная теорема.} \textit{Если $h\in Y$ и вектор $\vec h$ равен нулю, то функция $h$ есть тождественный нуль.}

Мы рекомендуем читателю доказать основную теорему самостоятельно или прочесть доказательство в приложении.

Из основной теоремы сразу следует формула (б). Действительно, из первой и пятой строк таблицы 1 мы получаем, что вектор, соответствующий функции $P-2S$, нулевой. Так как функция $P-2S$ удобная, то по основной теореме $P-2S=0$ для любой фигуры $F$.

\end{paracol}

\end{document}
